\section{The weighted quadratic row estimate}
\label{sec:ex-row-quadratic-price}

The part of the one-step action quadratic in the row field has a uniform
coefficient price. Its cost is the square of the exponentially weighted row
budget. The proof is an exact two-index expansion followed by a summable
Cauchy estimate; no finite-volume truncation of the row family is needed.

Let $\operatorname{CTax}$ be the iterated Cauchy derivative factor defined
in Section~\ref{sec:ex-wilson-vacuum-price}. Let $\mathcal A$ be the finite set
of plaquette types. Write $d_{\rm col}$ for
the number of color coordinates and $d_{\rm mat}$ for the dimension of the
matrix representation; these constants have different roles below. For each
$a\in\mathcal A$ and $i\in\mathbb N$, let $s_{a,i}(U;x)$ be a color vector
obtained by evaluating one path slot at the translated anchor $x$, and let
$\ell_{a,i}$ be that slot's length. Let $c_{a,i}\in\mathbb R$ be its row
coefficient, and assume $\ell_{a,i}\in\mathbb N$. The color vector takes values
in the finite-dimensional real Euclidean color space on physical fields; its
complex extension applies the complex-linear extension of the real color
projection to half the difference between the forward slot word and its
reversed word. On real fields this is the original skew-color slot. The
quadratic extension is bilinear, $\sum_c s_c s_c$, not a Hermitian norm.
Assume the row budget
\begin{equation}\label{eq:ex-row-budget}
 \sum_i |c_{a,i}|e^{\sigma\ell_{a,i}}\leq C
       \quad(a\in\mathcal A),\qquad \sigma>0.
\end{equation}

We use the following elementary path bound. For each $\sigma>0$ there are a
common complex coordinate radius $R_{\sigma/2}>0$ and a constant $S_G$ such
that, on the radius-$R_{\sigma/2}$ ball about the identity, every color
component of a slot of length $\ell$ has a holomorphic extension bounded by
$S_Ge^{(\sigma/2)\ell}$. More explicitly, if $H_a(z)$ is the matrix product along the slot walk and
$H_a^{\mathrm{rev}}(z)$ the product along the reversed slot walk, the
extension is
\[
 s_a^{\mathbb C}(z)=\Pi_{G,\mathbb C}\!\left(
       \frac{H_a(z)-H_a^{\mathrm{rev}}(z)}{2}\right),
\]
where $\Pi_{G,\mathbb C}$ is the complex-linear extension of the real
skew-color projection. On real links, the reversed full slot walk is the inverse walk; its string
matrix is the adjoint of the forward string matrix. Their half-difference is
skew-adjoint, and the physical color projection of this skew matrix is exactly
the transported-curvature slot. Thus the complex extension agrees with the
original real slot on physical fields. The two words each contain at most $2\ell$ edge
factors. Choose the radius so that each complexified edge matrix and its
inverse has operator norm at most $e^{\sigma/4}$; this is possible by
continuity near the compact unitary group. Each word then has norm at most
$e^{(\sigma/2)\ell}$. Apply the bounded complex-linear projection; its norm
and the fixed factor from the difference are absorbed in $S_G$. This proves
the bound uniformly in location and length.

Define the quadratic row action on a torus of side $m$ by
\[
 Q^{(2)}_m(U)=\frac12\sum_{x\in(\mathbb Z/m\mathbb Z)^4}
       \sum_{a\in\mathcal A}\left\|\sum_i c_{a,i}s_{a,i}(U;x)\right\|^2.
\]
The sum over $i$ converges absolutely on the real compact configuration
space by \eqref{eq:ex-row-budget} and the path bound. On the common complex
coordinate ball, the two-index expansion below also converges uniformly: its
absolute majorant is the product of the two row sums with exponent $\sigma/2$,
which is bounded by the assumed row budget. Thus the complexified action is
holomorphic there before any derivatives or momentum coefficients are taken.

\begin{proposition}[Uniform price of the quadratic row action]
\label{prop:ex-row-quadratic-price}
For each $t\in\mathbb N$, put $n=t+1$ and
\[
 \eta_t=\min\!\left(1,\frac{\sigma}{8n^2}\right).
\]
There is one coefficient picture $\mathcal Q_t$ valid at every torus side,
holomorphic for $t$ independent complex momenta of width $\eta_t$, whose
components are bounded by
\[
 \frac{|\mathcal A|d_{\rm col}}{2}C^2S_G^2
       \frac{\operatorname{CTax}(R_{\sigma/2},n)}{n!}.
\]
For each ordered list of $n$ Fourier-mode variations, its coefficient is zero
unless their total momentum vanishes; in the zero case it is the torus
cardinality times the value of $\mathcal Q_t$ at the independent momenta.
\end{proposition}

\begin{proof}
For each $a,i,j$, expand the squared norm into the two-slot term
\[
 \Gamma_{a,i,j}(U;x)=\frac12c_{a,i}c_{a,j}
       \sum_{c=1}^{d_{\rm col}}s_{a,i,c}(U;x)s_{a,j,c}(U;x).
\]
It has support length at most $\ell_{a,i}+\ell_{a,j}$ and total absolute
coefficient $d_{\rm col}|c_{a,i}c_{a,j}|/2$. Therefore its total weighted
activity satisfies, by Tonelli and \eqref{eq:ex-row-budget},
\[
 \begin{aligned}
 \sum_{a,i,j}\frac{d_{\rm col}}2|c_{a,i}c_{a,j}|
          e^{\sigma(\ell_{a,i}+\ell_{a,j})}
 &=\frac{d_{\rm col}}2\sum_a
       \left(\sum_i|c_{a,i}|e^{\sigma\ell_{a,i}}\right)^2\\
 &\leq \frac{|\mathcal A|d_{\rm col}}2C^2.
 \end{aligned}
\]

We now estimate the coefficient of each two-slot term. On the common
radius-$R_{\sigma/2}$ ball, the product of its two holomorphic slot extensions
is bounded by
$S_G^2e^{(\sigma/2)(\ell_{a,i}+\ell_{a,j})}$. The Cauchy derivative
estimate on the nested radii
$R_{\sigma/2}/2,\ldots,R_{\sigma/2}/2^n$ contributes
$\operatorname{CTax}(R_{\sigma/2},n)/n!$ to the normalized order-$n$
coefficient. Restricting a width-$\eta_t$ Fourier mode to a support of length
$\ell$ costs at most $e^{4n\eta_t\ell}$; for $n$ inputs the total additional
factor is at most $e^{4n^2\eta_t\ell}$. Since
$4n^2\eta_t\leq\sigma/2$, this factor is at most $e^{(\sigma/2)\ell}$.
Consequently the slot-extension and Fourier-mode factors together are bounded
by $S_G^2e^{\sigma(\ell_{a,i}+\ell_{a,j})}$.

The weighted activity sum just computed is thus a summable majorant for the
series of coefficient pictures, uniformly on the smaller momentum strip.
The Weierstrass $M$-test gives uniform convergence on compact subsets; its sum
is holomorphic, and Cauchy's integral formula permits termwise extraction of
each Taylor coefficient. On the real field, absolute convergence also permits
expanding the row square and interchanging the sums, so this summed picture
represents exactly $Q^{(2)}_m$.

Finally, translating each local two-slot term over all anchors multiplies its
Fourier coefficient by the character of the total momentum. The character sum
on the finite torus is zero at nonzero total momentum and equals the torus
cardinality at zero total momentum. For each pair atom, its finite support lies in a max-distance box of radius
$d=\ell_{a,i}+\ell_{a,j}$. Compute its local symbol on the auxiliary torus
$\Lambda_{2d+1}$; this box has no internal identifications there. Pull any
period-$m$ field back to this auxiliary torus by
$z(y)=z_m(\operatorname{wrap}_m(\operatorname{lift}y))$. The identity holds
even when $m<2d+1$, so repeated links are retained. Summing the local symbols
over atoms gives one picture independent of $m$, with the asserted volume
factor and price.
\end{proof}

\begin{corollary}[The averaged path-kernel coefficient price]
\label{cor:ex-row-controlled-action}
Suppose the unaveraged path kernel has the exact decomposition
\[
 Q_{p,R,m}=\lambda W_m+Q^{(2)}_m,\qquad 0\leq\lambda\leq Z,
\]
with $Q^{(2)}_m$ as above. Then the unaveraged order-$n$ picture has width
$\eta_t$ and price at most
\[
 A_{raw}=ZA_t+\frac{|\mathcal A|d_{\rm col}}{2}C^2S_G^2
       \frac{\operatorname{CTax}(R_{\sigma/2},n)}{n!},
       \qquad n=t+1.
\]
Let $\mathcal S$ be the finite signed-coordinate group and define the
average $Q^\sharp(U)=|\mathcal S|^{-1}\sum_{s\in\mathcal S}Q_{p,R,m}(\rho_sU)$,
where $\rho_s$ is the corresponding signed coordinate action on link
fluctuations. For this averaged kernel, the common picture has width $\eta_t$ and price
at most
\[
 A_{avg}=\operatorname{ReflectBill}(t,\eta_t)A_{raw},\qquad
 \operatorname{ReflectBill}(t,\eta)=\exp\bigl((t+1)^2\eta\bigr).
\]
Both coefficient identities hold at every torus side, with the exact
zero-total-momentum volume factor.
\end{corollary}

\begin{proof}
The unaveraged assertion follows by adding the bare-action picture and the
quadratic picture, after restricting the former to width $\eta_t\leq1$.
For completeness, the signed action on a picture entry is
\[
 (\operatorname{rot}_s L)(x,(j_r)_{r=1}^n,p)=
 \left(\prod_{r=1}^{n}\operatorname{reflectFee}_s
       (\operatorname{plotted}(p,r),k_r)\right)
 L\bigl(x,(k_r)_{r=1}^n,\operatorname{rotatePhase}_s(p)\bigr),
 \]
Here $n=t+1$ and $j_r\in\{0,1,2,3\}$ is the coordinate axis attached to the
$r$th slot and $k_r=s_{\rm perm}^{-1}(j_r)$ is the corresponding input axis,
used both in the fee and in the original picture's axis argument. Thus the
original picture is evaluated on the tuple $(k_1,\ldots,k_n)$. Each factor $\operatorname{reflectFee}_s$ is either $1$ or minus
the Fourier character $\exp(i\langle p,-e_{s(k_r)}\rangle)$ of one unit
direction. On the width-$\eta_t$ strip, each independent momentum has
imaginary coordinates bounded by $\eta_t$, and the plotted momentum
(including the dependent $p_0=-\sum_{r=1}^t p_r$) has imaginary coordinates
bounded by $n\eta_t$. Thus each such factor has price at most
$e^{n\eta_t}$. The product of the $n$ factors therefore costs at most
$e^{n^2\eta_t}$. The map $p\mapsto\operatorname{rotatePhase}_s(p)$ applies
the signed coordinate action $s^{-1}$ independently to each independent
momentum; the dependent momentum transforms by the same action because it is
minus their sum. This preserves the coordinate strip and incurs no additional
exponential price. This is precisely the picture transformation formula:
the character factors account for the translated base point in the signed
coordinate action, and the remaining coefficient is the original picture
evaluated on the reindexed directions and transformed momenta. Hence each
rotated picture has price at most
$\operatorname{ReflectBill}(t,\eta_t)A_{raw}$. The normalized finite signed
average preserves the same bound by the triangle inequality. Linearity of
the coefficient identity shows that the averaged picture represents the
averaged kernel and retains the same zero-momentum volume factor. This is the exact
orientation-average step needed for the input kernel.
\end{proof}
